% Môn: Vật lý
% Lớp: 12
% Chương: Chương I. Vật lí nhiệt
% Bài: L12C1 Bài 6. Nhiệt hóa hơi riêng · Bài toán xác định nhiệt nóng chảy riêng bằng thực nghiệm (thiết kế thí nghiệm với nước đá, đo đạc khối lượng, thời gian,
% ===== Câu 116 =====
% ID: L12C1B6-51-TLN
% Mức: VD
\dangbt{Bài toán xác định nhiệt nóng chảy riêng bằng thực nghiệm (thiết kế thí nghiệm với nước đá, đo đạc khối lượng, thời gian,}
\begin{ex}
% ID: L12C1B6-51-TLN
% Mức: VD
Trong quá trình chất khí truyền nhiệt và sinh công thì $A$ và $Q$ trong biểu thức $\Delta U=Q+A$ phải thỏa mãn điều kiện nào sau đây?
\shortans{D}
\loigiai{
• Chất khí truyền nhiệt và sinh công suy ra $Q<0, A<0$.
}
\end{ex}

% ===== Câu 197 =====
% ID: L12C1B6-231-TN
% Mức: TH
\begin{ex}
% ID: L12C1B6-231-TN
% Mức: TH
Trong thí nghiệm thực hành xác định nhiệt nóng chảy riêng của nước đá bằng nhiệt lượng kế, tại sao người ta cần sử dụng nước đá đang tan, tức đang ở đúng $0^{\circ}\mathrm{C}$, mà không sử dụng nước đá lấy trực tiếp từ ngăn đông sâu đang ở nhiệt độ âm?
\choice
{Để tránh làm hỏng nhiệt kế điện tử và thành bình nhiệt lượng kế khi tiếp xúc với nhiệt độ quá lạnh}
{Để nước đá dễ dàng tan nhanh hơn, giúp rút ngắn thời gian làm thí nghiệm}
{\True Để nhiệt lượng do nguồn điện cung cấp chỉ dùng cho quá trình nóng chảy của nước đá mà không tốn thêm năng lượng để tăng nhiệt độ của đá từ nhiệt độ âm lên $0^{\circ}\mathrm{C}$}
{Để khối lượng nước đá không bị thay đổi trong suốt quá trình chuẩn bị cân đo khối lượng}
\loigiai{
Nếu sử dụng nước đá ở nhiệt độ âm, một phần nhiệt lượng cung cấp sẽ dùng để làm nước đá tăng nhiệt độ đến $0^{\circ}\mathrm{C}$ trước khi nóng chảy. Khi dùng nước đá đang tan, toàn bộ nhiệt lượng cung cấp chủ yếu dùng cho quá trình chuyển thể từ rắn sang lỏng ở nhiệt độ không đổi $0^{\circ}\mathrm{C}$, giúp phép tính $Q=\lambda m$ đơn giản và chính xác hơn.
}
\end{ex}

% ===== Câu 199 =====
% ID: L12C1B6-232-TN
% Mức: TH
\begin{ex}
% ID: L12C1B6-232-TN
% Mức: TH
Một nhóm học sinh thiết kế phương án thí nghiệm đo nhiệt nóng chảy riêng của nước đá. Họ sử dụng một dây điện trở đặt trong phễu chứa nước đá đang tan, bên dưới phễu đặt một chiếc cốc trên cân điện tử để đo khối lượng nước chảy ra. Để xác định lượng nước đá tan do nhiệt lượng của dây điện trở cung cấp và loại trừ sự tan chảy tự nhiên do môi trường, nhóm học sinh cần thực hiện phép đo khối lượng như thế nào?
\choice
{Chỉ cần đo khối lượng nước trong cốc sau một khoảng thời gian đun $t$ từ khi bật nguồn điện}
{Lấy khối lượng nước đá ban đầu trong phễu trừ đi khối lượng nước đá còn lại sau thời gian $t$}
{\True Đo khối lượng nước chảy vào cốc sau cùng một khoảng thời gian $t$ trong hai trường hợp: khi chưa bật nguồn điện $(m_1)$ và khi đã bật nguồn điện $(m_2)$, từ đó khối lượng nước đá tan do nguồn điện cung cấp là $m_2-m_1$}
{Đo khối lượng nước chảy vào cốc khi chưa bật nguồn điện, rồi cộng thêm khối lượng nước đá còn lại trong phễu}
\loigiai{
Nước đá trong phễu luôn có thể tan tự nhiên do nhận nhiệt từ môi trường. Vì vậy, cần đo khối lượng nước đá tan tự nhiên $m_1$ trong thời gian $t$ trước khi bật nguồn điện, sau đó đo khối lượng nước đá tan $m_2$ trong cùng thời gian $t$ khi bật nguồn điện. Khối lượng nước đá tan do điện là
 $m=m_2-m_1.$
}
\end{ex}

% ===== Câu 201 =====
% ID: L12C1B6-233-TN
% Mức: TH
\begin{ex}
% ID: L12C1B6-233-TN
% Mức: TH
Khi đo nhiệt nóng chảy riêng của nước đá bằng bộ dụng cụ thí nghiệm thực hành, nếu học sinh bỏ qua nhiệt lượng hao phí truyền ra môi trường xung quanh và truyền cho các thiết bị đo trong quá trình tính toán theo công thức $\lambda=\dfrac{Pt}{m}$ thì kết quả thu được sẽ như thế nào so với giá trị thực tế?
\choice
{Nhỏ hơn giá trị thực tế vì khối lượng nước đá tan $m$ thực tế lớn hơn}
{\True Lớn hơn giá trị thực tế vì một phần năng lượng thực tế đã bị hao phí ra môi trường, làm cho lượng nước đá nóng chảy thực tế nhỏ hơn lượng đáng lẽ phải tan nếu không có hao phí}
{Bằng giá trị thực tế vì hao phí nhiệt lượng ra môi trường là không đáng kể}
{Nhỏ hơn giá trị thực tế vì công suất thực tế cung cấp cho nước đá bị giảm đi}
\loigiai{
Trong thực tế, một phần nhiệt lượng do dây nung tỏa ra bị thất thoát ra môi trường và làm nóng dụng cụ thí nghiệm. Do đó, nhiệt lượng thực sự truyền cho nước đá nhỏ hơn $Pt$. Khi vẫn lấy $\lambda=\dfrac{Pt}{m}$, giá trị $\lambda$ tính được sẽ lớn hơn giá trị thực tế.
}
\end{ex}

% ===== Câu 203 =====
% ID: L12C1B6-234-TN
% Mức: VD
\begin{ex}
% ID: L12C1B6-234-TN
% Mức: VD
Tiến hành thí nghiệm đo nhiệt nóng chảy riêng của nước đá bằng một dây nung có công suất tỏa nhiệt không đổi $P=20\,\mathrm{W}$ đặt trong phễu chứa nước đá đang tan. Trong thời gian $t=3$ phút kể từ khi bật nguồn điện, khối lượng nước chảy vào cốc tăng thêm $11\,\mathrm{g}$ so với khi chưa bật nguồn điện. Tính nhiệt nóng chảy riêng của nước đá xác định từ kết quả thực nghiệm này.
\choice
{$3,60\cdot 10^5\,\mathrm{J/kg}$}
{$3,40\cdot 10^5\,\mathrm{J/kg}$}
{\True $3,27\cdot 10^5\,\mathrm{J/kg}$}
{$3,00\cdot 10^5\,\mathrm{J/kg}$}
\loigiai{
Đổi $t=3$ phút $=180\,\mathrm{s}$ và $m=11\,\mathrm{g}=0,011\,\mathrm{kg}$. Nhiệt lượng do dây nung cung cấp:
 $Q=Pt=20\cdot 180=3600 \mathrm{J}.$
 Nhiệt nóng chảy riêng của nước đá:
 $\lambda=\dfrac{Q}{m}=\dfrac{3600}{0,011}\approx 3,27\cdot 10^5 \mathrm{J/kg}.$
}
\end{ex}

% ===== Câu 205 =====
% ID: L12C1B6-235-TN
% Mức: VD
\begin{ex}
% ID: L12C1B6-235-TN
% Mức: VD
Trong một bài thực hành đo nhiệt nóng chảy riêng của nước đá, một nhóm học sinh sử dụng bếp điện có công suất phát nhiệt có ích thực tế truyền cho nước đá nóng chảy là $P=50\,\mathrm{W}$. Đồ thị thực nghiệm biểu diễn khối lượng nước chảy ra theo thời gian đun là một đường thẳng. Họ xác định được trong khoảng thời gian $\Delta t=25$ phút, lượng nước đá nóng chảy tăng thêm $m=180\,\mathrm{g}$ so với khi chưa bật điện. Hãy tính nhiệt nóng chảy riêng của nước đá xác định từ kết quả thực nghiệm này.
\choice
{$3,34\cdot 10^5\,\mathrm{J/kg}$}
{$3,40\cdot 10^5\,\mathrm{J/kg}$}
{\True $4,17\cdot 10^5\,\mathrm{J/kg}$}
{$3,15\cdot 10^5\,\mathrm{J/kg}$}
\loigiai{
Đổi $\Delta t=25$ phút $=1500\,\mathrm{s}$ và $m=180\,\mathrm{g}=0,18\,\mathrm{kg}$. Nhiệt lượng có ích cung cấp cho nước đá nóng chảy là
 $Q=P\Delta t=50\cdot 1500=75000 \mathrm{J}.$
 Nhiệt nóng chảy riêng đo được từ thực nghiệm:
 $\lambda=\dfrac{Q}{m}=\dfrac{75000}{0,18}\approx 4,17\cdot 10^5 \mathrm{J/kg}.$
}
\end{ex}

% ===== Câu 207 =====
% ID: L12C1B6-236-TN
% Mức: VDC
\begin{ex}
% ID: L12C1B6-236-TN
% Mức: VDC
Một học sinh tiến hành thí nghiệm xác định nhiệt nóng chảy riêng của nước đá bằng nhiệt lượng kế có dây nung với công suất $\mathcal{P} = 24\ \mathrm{W}$, cùng với cân và cốc hứng nước. Sơ đồ thí nghiệm như hình bên. 
 • Bước 1: Cho nước đá tan ở $0^\circ \mathrm{C}$ vào nhiệt lượng kế sau thời gian $t = 360\ \mathrm{s}$ thì thu được khối lượng nước ở cốc là $m = 4 \times 10^{-3}\ \mathrm{kg}$.
 
 • Bước 2: Bật biến áp nguồn để nung nóng lượng nước đá cũng trong thời gian $360\ \mathrm{s}$. Sau đó học sinh ghi nhận tổng lượng nước trong cốc là $M = 34 \times 10^{-3}\ \mathrm{kg}$.
 
 Bước 2 được tiến hành ngay sau bước 1, xem dây nung không tỏa nhiệt ra môi trường. Nhiệt nóng chảy riêng của nước đá mà học sinh này đo được xấp xỉ.
\choice
{$2,54.10^5\ J/kg$}
{\True $3,32.10^5\ J/kg$}
{$2,88.10^5\ J/kg$}
{$6,65.10^5\ J/kg$}
\loigiai{
Công thức tính nhiệt nóng chảy riêng: $L = \frac{Q}{m} = \frac{P t}{\Delta m}$, với $P$ công suất, $t$ thời gian, $\Delta m$ khối lượng nước đá tan do dây nung. 
 **A.** Sai vì tính toán nhiệt nóng chảy riêng thấp hơn giá trị đúng, có thể do sai khối lượng nước đá tan do dây nung. 
 **B.** Đúng vì tính đúng khối lượng nước đá tan do dây nung là $26\,g$, công suất và thời gian đúng, kết quả phù hợp với giá trị chuẩn của nước đá. 
 **C.** Sai vì giá trị nhiệt nóng chảy riêng thấp hơn đáp án đúng, do sai số trong xác định khối lượng nước đá tan. 
 **D.** Sai vì giá trị quá cao, không phù hợp với thực tế và dữ liệu thí nghiệm. 
 Lưu ý: Trong lời giải có lỗi ghi $3,32.10^{-5}$ thay vì $3,32.10^{5}$ J/kg, cần sửa lại cho đúng đơn vị và giá trị.
}
\end{ex}
